\documentclass[11pt]{article}
\usepackage{amsmath,amssymb,amsthm,hyperref}
\usepackage[margin=1in]{geometry}
\newtheorem{theorem}{Theorem}
\title{Every prescribed least period in a complex rational delay recurrence}
\author{Principia Math In Progress}
\date{29 September 2026}
\begin{document}\maketitle
\begin{abstract}
For every integer $p\ge2$, we construct an everywhere-defined solution of
least period $p$ of $z_{n+1}=(\alpha+\beta z_{n-k})/(\gamma-z_n)$, with
$k=p-1$ and all three parameters nonzero and nonreal. The construction
proves the higher-period existence conjecture in MathDB 333488.
\end{abstract}

\section*{Task 1: a full construction for every period}
Conjecture 3.1 on page 14 of Hassan's manuscript \cite{Hassan} asks for
arbitrarily high periods of this family of delay recurrences. Pages 1--2
allow arbitrary complex parameters and initial conditions, and the
examples immediately preceding the conjecture vary the delay and
parameters. No fixed-delay or fixed-parameter restriction is stated.
We prove the precise stronger existence statement that every prescribed
integer $p\ge2$ occurs as a least period, with explicit parameters,
delay, and initial history. This is an existence result in the family;
it makes no claim that one fixed recurrence realizes all periods.

\begin{theorem}
Fix an integer $p\ge2$ and put
\[
 \theta=\frac{\pi}{p},\qquad t=2\cos\theta,\qquad s=1+i,
 \qquad k=p-1,
\]
\[
 \alpha=s^2=2i,\qquad \beta=s=1+i,\qquad
 \gamma=s(1+t).
\]
Let $T(w)=1/(t-w)$, set $w_0=2i$, and define $w_j=T^j(w_0)$ for
$j\ge0$. Then $w_j$ is periodic with least period $p$. Extend it
periodically to all integer indices and set $z_j=s w_j$.
The initial history $z_{-k},\ldots,z_0$ produces an everywhere-defined
solution of
\[
 z_{n+1}=\frac{\alpha+\beta z_{n-k}}{\gamma-z_n}
 \qquad(n\ge0)
\]
of least period $p$. All three parameters are nonzero and nonreal.
\end{theorem}
\begin{proof}
First, every $w_j$ lies in the upper half-plane. Indeed, $t$ is real and
\[
 \operatorname{Im}T(w)=\frac{\operatorname{Im}w}{|t-w|^2}>0
 \qquad(\operatorname{Im}w>0).
\]
Thus every iterate is finite and the denominator $t-w_j$ never vanishes.

Set $u=e^{i\theta}$ and $v=e^{-i\theta}$. These are distinct, satisfy
$u+v=t$ and $uv=1$, and are the fixed points of $T$. Direct subtraction
gives
\[
 T(w)-u=\frac{u(w-u)}{t-w},\qquad
 T(w)-v=\frac{v(w-v)}{t-w}.
\]
Consequently the injective fractional linear map
$R(w)=(w-u)/(w-v)$ satisfies
\[
 R(T(w))=e^{2i\theta}R(w).
\]
This identity is defined throughout the orbit because $v$ lies in the
lower half-plane. Also $R(w_0)\ne0$, since $|w_0|=2$ and $|u|=1$.
It follows by induction that
\[
 R(w_j)=e^{2\pi i j/p}R(w_0).
\]
The multiplier is a primitive $p$th root of unity, so injectivity of
$R$ gives $w_j=w_0$ if and only if $p$ divides $j$. In particular,
$w_p=w_0$, and iteration gives $w_{j+p}=w_j$ for all $j\ge0$.
This proves periodicity and minimality, not just an upper bound on
the period. The periodic extension also satisfies $w_{j+1}=T(w_j)$
for all integers $j$.

Now $k=p-1$ implies $w_{n-k}=w_{n+1}$. The defining equation of $T$
is $w_{n+1}(t-w_n)=1$, and hence
\[
 1+w_{n-k}=1+w_{n+1}=(1+t-w_n)w_{n+1}.
\]
Because $w_n$ is nonreal, the factor $1+t-w_n$ is nonzero.
Substituting the stated parameters therefore gives
\[
 \frac{\alpha+\beta z_{n-k}}{\gamma-z_n}
 =\frac{s^2+s^2w_{n-k}}{s(1+t-w_n)}
 =s w_{n+1}=z_{n+1}.
\]
No pole is introduced, since $s\ne0$. Multiplication by $s$ preserves
least period. Finally $0\le t<2$, so $1+t>0$ and the imaginary parts
of $\alpha,\beta,\gamma$ are respectively $2,1,1+t$, all positive.
\end{proof}

The use of $p$ as the least period agrees with the source's ``prime
period'' definition. In particular the result supplies arbitrarily
large periods without counting multiples of a shorter period. The
construction does not assert attraction or chaos.

\section*{Four proposed follow-up tasks}
\begin{enumerate}
\item[Task 2.] Classify the parameters for which the ansatz
$z_{n-k}=z_{n+1}$ reduces the recurrence to a finite-order fractional
linear transformation, retaining all pole exclusions.
\item[Task 3.] For each prescribed period, determine the smallest
positive delay that admits a solution of that least period.
\item[Task 4.] Determine the stability of these periodic solutions
under perturbations of the full initial history of length $k+1$.
\item[Task 5.] Decide whether a single fixed parameter triple and
fixed delay can admit periodic solutions of unbounded least period.
\end{enumerate}

\begin{thebibliography}{9}
\bibitem{Hassan} Sk.~Sarif Hassan,
\emph{Dynamics of the complex rational delay recursive sequence
$z_{n+1}=(\alpha+\beta z_{n-k})/(\gamma-z_n)$},
\href{https://arxiv.org/abs/1606.08887v1}{arXiv:1606.08887v1} (2016),
pp.~1--2 and Conjecture~3.1, p.~14.
\end{thebibliography}
\end{document}
